\section{问题五模型的建立与求解}

\subsection{本问求解思路}

问题五要求舞龙队沿问题四已经确定的路径行进，龙头前把手速度保持不变，并求出龙头前把手的最大允许速度，使得全队所有把手速度均不超过 $2\,\mathrm{m/s}$。根据 2024A 项目模型边界规则，本问直接沿用问题四中的统一路径模型和速度递推关系，不再重复建立调头路径几何、S 形圆弧构造、位置递推过程和速度递推来源，仅新增单位龙头速度下的速度放大系数、全路径最大速度放大系数以及最大龙头速度计算模型。

问题四已经在龙头前把手速度为 $1\,\mathrm{m/s}$ 的条件下，给出了沿统一路径 $\Gamma(s)$ 的全队位置与速度求解方法。由于路径几何形状固定，且刚性速度投影关系对速度是线性的，若将龙头前把手速度由 $1\,\mathrm{m/s}$ 改为某一常数 $v_0$，则任一给定路径构型下所有把手的速度均随 $v_0$ 成比例缩放。因此，本问先在单位龙头速度下计算全路径、全把手的最大速度放大系数，再由速度上限 $2\,\mathrm{m/s}$ 反推出龙头前把手最大速度。

为展示问题五的建模与求解流程，绘制求解流程图，如图~\ref{fig:p5-flow} 所示。

\begin{figure}[H]
    \centering
    \includegraphics[width=0.82\textwidth]{figures/p5_flow.png}
    \caption{问题五求解流程图}
    \label{fig:p5-flow}
\end{figure}

图~\ref{fig:p5-flow} 给出了本问从单位速度场计算到最大速度放大系数搜索，再到龙头最大速度确定的主要过程。

\subsection{单位龙头速度下的速度放大系数}

设龙头前把手对应的路径弧长参数为 $s_0$。在单位龙头速度下，沿用问题四速度递推关系，令
\begin{equation}
    q_0(s_0)=1.
    \label{eq:p5-q0}
\end{equation}
对给定 $s_0$，调用问题四的位置递推模型得到全部把手弧长参数 $s_i(s_0)$ 与位置 $P_i(s_0)$，再调用问题四速度递推式~\eqref{eq:p4-speed-recursion} 得到
\begin{equation}
    q_i(s_0),\qquad i=0,1,\ldots,223.
    \label{eq:p5-qi}
\end{equation}
由于问题四路径采用弧长参数化，单位龙头速度下第 $i$ 个把手的速度大小即为
\begin{equation}
    \lambda_i(s_0)=|q_i(s_0)|.
    \label{eq:p5-lambda}
\end{equation}
其中 $\lambda_i(s_0)$ 表示第 $i$ 个把手相对于龙头前把手的速度放大系数。当 $\lambda_i(s_0)>1$ 时，表示该把手速度大于龙头速度；当 $\lambda_i(s_0)<1$ 时，表示该把手速度小于龙头速度。

\subsection{最大速度优化模型}

问题五考察问题四设定路径上 $t\in[-100,100]$ 的运动过程。在单位龙头速度下有 $s_0=t$，故可将搜索区间写为
\begin{equation}
    s_0\in[-100,100].
    \label{eq:p5-range}
\end{equation}
定义全路径、全把手的最大速度放大系数为
\begin{equation}
    \Lambda_{\max}
    =
    \max_{s_0\in[-100,100]}
    \max_{0\le i\le223}
    \lambda_i(s_0).
    \label{eq:p5-max-lambda}
\end{equation}

当龙头前把手速度调整为常数 $v_0$ 时，路径构型不变，只是时间尺度发生线性变化。由速度递推关系的线性性，第 $i$ 个把手实际速度为
\begin{equation}
    v_i(s_0;v_0)=v_0\lambda_i(s_0).
    \label{eq:p5-actual-speed}
\end{equation}
题目要求所有把手速度均不超过 $2\,\mathrm{m/s}$，因此必须满足
\begin{equation}
    v_0\lambda_i(s_0)\le2,
    \qquad
    s_0\in[-100,100],\quad i=0,1,\ldots,223.
    \label{eq:p5-speed-limit}
\end{equation}
由式~\eqref{eq:p5-max-lambda} 和式~\eqref{eq:p5-speed-limit} 可得龙头前把手最大允许速度为
\begin{equation}
    v_{0,\max}=\frac{2}{\Lambda_{\max}}.
    \label{eq:p5-vmax}
\end{equation}
因此，问题五被转化为单位龙头速度下的最大速度放大系数搜索问题。

\subsection{数值求解方法}

由于 $\Lambda_{\max}$ 由全路径和全部把手共同决定，本文采用“全区间粗搜索、活跃区间加密、局部极值精修”的方法求解。首先在 $[-100,100]$ 内逐步扫描龙头弧长位置 $s_0$，对每个 $s_0$ 调用问题四的位置与速度递推模型，计算当前构型下的
\[
    \Lambda(s_0)=\max_{0\le i\le223}\lambda_i(s_0).
\]
然后找到 $\Lambda(s_0)$ 较大的活跃区间，并在该区间内进行加密搜索，直到 $\Lambda_{\max}$ 的数值稳定。

\begin{algorithm}[H]
    \caption{问题五龙头最大速度求解算法}
    \label{alg:p5-vmax}
    \begin{algorithmic}[1]
        \Require 问题四统一路径与速度递推模型，速度上限 $2\,\mathrm{m/s}$
        \Ensure 龙头前把手最大允许速度 $v_{0,\max}$
        \State 令 $\Lambda_{\max}=1$
        \For{$s_0$ 在 $[-100,100]$ 内粗搜索}
            \State 调用问题四位置递推模型得到全部 $s_i(s_0)$ 与 $P_i(s_0)$
            \State 令 $q_0(s_0)=1$，调用问题四速度递推模型得到全部 $q_i(s_0)$
            \State 计算 $\Lambda(s_0)=\max_i |q_i(s_0)|$
            \State 若 $\Lambda(s_0)>\Lambda_{\max}$，则更新 $\Lambda_{\max}$
        \EndFor
        \State 对最大值附近活跃区间进行加密搜索，得到最终 $\Lambda_{\max}$
        \State 输出 $v_{0,\max}=2/\Lambda_{\max}$
    \end{algorithmic}
\end{algorithm}

\subsection{求解结果与分析}

按照上述模型计算得到，全路径、全把手最大速度放大系数为
\begin{equation}
    \Lambda_{\max}=1.604793379.
    \label{eq:p5-lambda-result}
\end{equation}
该最大值出现在单位龙头速度计时下
\begin{equation}
    t_c=14.479969\,\mathrm{s}
    \label{eq:p5-critical-time}
\end{equation}
附近。此时龙头前把手已经完成 S 形调头并进入盘出段，而龙头后方部分把手仍位于调头曲线附近，路径曲率差异使速度传递出现最大放大。

由式~\eqref{eq:p5-vmax} 得龙头前把手最大允许速度为
\begin{equation}
    v_{0,\max}
    =\frac{2}{1.604793379}
    =1.246266358\,\mathrm{m/s}.
    \label{eq:p5-vmax-result}
\end{equation}
保留 $6$ 位小数，得到
\begin{equation}
    v_{0,\max}=1.246266\,\mathrm{m/s}.
    \label{eq:p5-vmax-six}
\end{equation}

临界状态下部分把手速度如表~\ref{tab:p5-critical-speed} 所示。

\begin{table}[H]
    \centering
    \caption{问题五临界状态下部分把手速度}
    \label{tab:p5-critical-speed}
    \begin{tabular}{c c c}
        \toprule
        把手 & 速度放大系数 $\lambda_i$ & 当 $v_0=1.246266358\,\mathrm{m/s}$ 时的速度 $v_i/(\mathrm{m/s})$ \\
        \midrule
        龙头前把手 & 1.000000 & 1.246266 \\
        第 $1$ 节龙身前把手 & 1.554248 & 1.937007 \\
        第 $2$ 节龙身前把手 & 1.554248 & 1.937007 \\
        第 $3$ 节龙身前把手 & 1.604793 & 2.000000 \\
        第 $4$ 节龙身前把手 & 1.604793 & 2.000000 \\
        第 $5$ 节龙身前把手 & 1.604793 & 2.000000 \\
        第 $6$ 节龙身前把手 & 1.604793 & 2.000000 \\
        第 $7$ 节龙身前把手 & 1.604793 & 2.000000 \\
        \bottomrule
    \end{tabular}
\end{table}

表~\ref{tab:p5-critical-speed} 表明，限制龙头速度的不是龙头前把手本身，而是调头曲线附近的龙身把手。当龙头速度取 $1.246266358\,\mathrm{m/s}$ 时，第 $3$ 至第 $7$ 节龙身前把手附近达到速度上限 $2\,\mathrm{m/s}$，因此该速度即为满足全队速度约束的理论最大值。

\subsection{可靠性检验}

\subsubsection{速度线性缩放检验}

对不同龙头速度 $v_0$ 代入式~\eqref{eq:p5-actual-speed}，可验证所有把手速度均与 $v_0$ 成正比例变化。由于位置构型只由 $s_0$ 决定，速度比例系数 $\lambda_i(s_0)$ 不随 $v_0$ 改变，因此单位速度场可用于反推任意常速龙头条件下的全队速度。

\subsubsection{临界性检验}

当
\[
    v_0=1.246266358\,\mathrm{m/s}
\]
时，有
\[
    v_0\Lambda_{\max}=2.000000\,\mathrm{m/s}.
\]
若将龙头速度增加 $10^{-6}\,\mathrm{m/s}$，则最大把手速度变为
\[
    (v_0+10^{-6})\Lambda_{\max}=2.000001605\,\mathrm{m/s}>2\,\mathrm{m/s},
\]
不满足题设速度上限；若将龙头速度减少 $10^{-6}\,\mathrm{m/s}$，则所有把手速度均小于 $2\,\mathrm{m/s}$。因此式~\eqref{eq:p5-vmax-result} 给出的速度具有临界性。

\subsubsection{搜索稳定性检验}

对 $[-100,100]$ 的搜索区间先采用粗步长定位最大值附近区间，再在活跃区间内加密搜索。加密后最大速度放大系数稳定为
\[
    \Lambda_{\max}=1.604793379,
\]
对应龙头最大速度稳定为
\[
    v_{0,\max}=1.246266358\,\mathrm{m/s}.
\]
结果在 $10^{-6}\,\mathrm{m/s}$ 精度下保持不变。

\subsection{模型评价}

本问将问题五转化为单位龙头速度下的速度放大系数搜索问题，避免了对不同龙头速度重复求解全队运动。模型利用了刚性速度递推的线性结构，计算效率较高，并能够指出限制龙头最大速度的具体位置。模型不足在于仍基于二维刚性平面运动假设，未考虑实际舞龙人员步态差异、操作误差和板凳连接间隙，因此所得最大速度是题设理想模型下的理论上限。实际表演中可在 $1.246266\,\mathrm{m/s}$ 的基础上留出安全裕度。
